\chapter{Lezione 1a - }

\section{Fisica dei Plasmi e Teoria della Dinamo: Lo Stato di Plasma}

Gli appunti che seguono sono tratti dal Modulo 1a del corso di Fisica dei Plasmi e Teoria della Dinamo (tenuto dal Prof. Alfio Bonanno, afferente all'INAF – Osservatorio Astrofisico di Catania). Il nucleo tematico di questa sezione introduttiva riguarda le caratteristiche fondamentali dello stato di plasma, con un'attenzione particolare rivolta ai fenomeni di schermo elettrostatico ("screening") e all'insorgenza delle oscillazioni collettive.

\subsection{Struttura e Bibliografia del Modulo}

Il programma di questa sezione si articola logicamente in tre macro-argomenti principali:
\begin{enumerate}
\item La presenza e le caratteristiche dei plasmi nell'Universo e negli ambienti di laboratorio.
\item Il fenomeno dello schermo di Debye e il principio di quasi-neutralità.
\item La natura e la dinamica delle oscillazioni di plasma.
\end{enumerate}

Per l'approfondimento teorico di tali concetti, si fa riferimento alle seguenti letture consigliate: il capitolo 1 del testo di Boyd \& Sanderson, il capitolo 1 del testo di Bittencourt e il capitolo 1 del testo di Fitzpatrick. Ulteriore materiale di riferimento include le note specifiche del corso (capitolo 1, paragrafi da 1.1 a 1.3) e il primo set di problemi assegnati per la pratica.

\section{I Plasmi nell'Universo e nel Laboratorio}

\subsection{Definizione Formale di Plasma}

Il punto di partenza concettuale per l'intero corso è rispondere a una domanda apparentemente semplice: che cos'è un plasma? In termini fisici rigorosi, un plasma viene definito come un gas \textbf{quasi-neutro} composto da particelle dotate di carica elettrica (ed eventualmente anche da una frazione di particelle neutre), il quale manifesta un \textbf{comportamento collettivo}.

Da un punto di vista storico, l'introduzione in fisica del termine "plasma" si deve a Irving Langmuir nel 1928. Osservando i gas ionizzati all'interno dei tubi a scarica elettrica, Langmuir notò un'analogia comportamentale con il plasma sanguigno: così come quest'ultimo trasporta corpuscoli biologici, il gas ionizzato fungeva da mezzo di trasporto fluido per ioni ed elettroni liberi.

\subsection{Le Tre Condizioni di Esistenza dello Stato di Plasma}

Affinché un generico gas ionizzato possa essere formalmente classificato e trattato matematicamente come un plasma, devono essere soddisfatte tre condizioni fisiche fondamentali. Tali condizioni, che verranno derivate e discusse in dettaglio nel corso della trattazione, sono le seguenti:

\begin{itemize}
\item \textbf{Quasi-neutralità macroscopica:} La prima condizione richiede che la lunghezza di Debye (indicata con $\lambda_D$) sia trascurabile rispetto alla scala di lunghezza caratteristica (o dimensione macroscopica) del sistema fisico, denotata con $L$. In termini matematici, ciò si esprime come:
\begin{equation}
\lambda_D \ll L
\end{equation}
Quando questa disuguaglianza è verificata, le cariche elettriche all'interno del mezzo vengono efficacemente schermate su distanze brevi, garantendo che, su scala macroscopica, il sistema mantenga una sostanziale neutralità elettrica (la cosiddetta "quasi-neutralità").


\item \textbf{Accoppiamento debole (Dominanza del comportamento collettivo):} La seconda condizione coinvolge il parametro $N_D$, ovvero il numero totale di particelle cariche contenute all'interno di una "sfera di Debye"[cite: 4]. Si richiede che tale parametro sia molto maggiore dell'unità:
\begin{equation}
    N_D = \frac{4\pi}{3}n\lambda_D^3 \gg 1
\end{equation}
dove $n$ rappresenta la densità numerica delle particelle[cite: 4]. La presenza di moltissime particelle all'interno del raggio d'azione schermante assicura che il sistema si trovi in un regime di accoppiamento debole ("weak coupling")[cite: 4]. In questo regime, l'interazione a lungo raggio tra miriadi di cariche domina nettamente rispetto alle collisioni binarie a corto raggio[cite: 4].

\item \textbf{Dominanza della dinamica del plasma sulle collisioni neutre:} La terza e ultima condizione riguarda l'evoluzione temporale del sistema, ponendo un vincolo sul prodotto tra la frequenza angolare tipica delle oscillazioni di plasma ($\omega_p$) e il tempo caratteristico di collisione con le particelle neutre ($\tau_n$)[cite: 4]. Questa relazione deve soddisfare:
\begin{equation}
    \omega_p \tau_n \gg 1
\end{equation}
Fisicamente, questa disequazione impone che la dinamica oscillatoria collettiva delle cariche elettriche avvenga su scale temporali nettamente più veloci rispetto a quelle necessarie per subire una collisione dissipativa con una particella neutra del gas di fondo[cite: 4]. Solo così il gas può mostrare un comportamento governato primariamente dalle forze elettromagnetiche.

\end{itemize}

\section{L'Onnipresenza dei Plasmi e la Diversità degli Ambienti Fisici}

Quasi tutta la materia barionica visibile nell'Universo si trova allo stato ionizzato, ed è dunque classificabile come plasma. Questo stato della materia è letteralmente onnipresente e si manifesta in una vastissima gamma di contesti astrofisici, tra cui gli interni e le atmosfere stellari, il mezzo interstellare, il mezzo intra-ammasso delle galassie (intracluster medium), i venti stellari, le magnetosfere planetarie o compatte e i dischi di accrescimento.

Alcuni esempi notevoli del manifestarsi dei plasmi in natura includono:
\begin{itemize}
\item La \textbf{corona solare}, costituita da un plasma estremamente caldo che raggiunge temperature dell'ordine di $10^6 \text{ K}$. Essa è osservabile dettagliatamente nella banda dell'estremo ultravioletto (EUV), ad esempio isolando l'emissione dello ione Fe IX alla lunghezza d'onda di $171 \text{ \AA}$ (come documentato dalle osservazioni del Solar Dynamics Observatory, NASA/SDO).
\item Le \textbf{aurore polari}, che sono fenomeni ottici generati nella ionosfera terrestre, caratterizzata da un plasma a temperatura nettamente inferiore, pari a circa $10^3 \text{ K}$. (Visibili chiaramente anche dallo spazio, come ripreso dalla Stazione Spaziale Internazionale durante la Spedizione 23).
\item L'\textbf{ammasso di galassie di Perseo}, il cui mezzo permeante è un plasma intergalattico a ben $10^8 \text{ K}$, rivelabile attraverso l'intensa emissione di radiazione di frenamento (bremsstrahlung) nei raggi X (ad esempio tramite le indagini dell'Osservatorio a Raggi X Chandra, NASA/CXC).
\end{itemize}

I plasmi conosciuti coprono un intervallo straordinario, estendendosi per circa $30$ ordini di grandezza in termini di densità numerica elettronica ($n_e$). I valori tipici di ordine di grandezza associati ai principali ambienti di plasma sono i seguenti:
\begin{itemize}
\item \textbf{Mezzo intra-ammasso (Intracluster medium):} $n_e \approx 10^{-3} \text{ cm}^{-3}$, temperatura $T \approx 10^{7.5} \text{ K}$ e intensità del campo magnetico $B \approx 10^{-6} \text{ G}$ (caratterizzato dall'emissione di bremsstrahlung nei raggi X).
\item \textbf{Mezzo interstellare ionizzato caldo (Warm ionised ISM):} $n_e \approx 10^{-1} \text{ cm}^{-3}$, $T \approx 10^4 \text{ K}$, $B \approx 5 \times 10^{-6} \text{ G}$ (analizzabile mediante le indagini sulla misura di dispersione, DM, e misura di rotazione, RM, del segnale delle pulsar).
\item \textbf{Vento solare a 1 UA:} $n_e \approx 5 \text{ cm}^{-3}$, $T \approx 10^5 \text{ K}$, $B \approx 5 \times 10^{-5} \text{ G}$ (descrivibile teoricamente come un plasma non collisionale).
\item \textbf{Ionosfera terrestre (Strato F):} $n_e \approx 10^6 \text{ cm}^{-3}$, $T \approx 10^3 \text{ K}$, $B \approx 0.5 \text{ G}$ (strato fondamentale in quanto responsabile della riflessione delle onde radio).
\item \textbf{Corona solare:} $n_e \approx 10^9 \text{ cm}^{-3}$, $T \approx 10^6 \text{ K}$, $B$ compreso tra $1$ e $100 \text{ G}$ (ambiente dinamico che costituisce il teatro in cui avvengono i brillamenti o "flares" e i processi di riconnessione magnetica).
\item \textbf{Magnetosfera delle pulsar:} $n_e \sim 10^{11} \text{ cm}^{-3}$, assenza di una chiara temperatura termica associabile al bulk, e $B \approx 10^{12} \text{ G}$ (dominata dalla fisica quantistica e relativistica a causa della presenza di plasmi formati da coppie elettrone-positrone, $e^{\pm}$).
\item \textbf{Fotosfera solare:} $n_e \approx 10^{13} \text{ cm}^{-3}$, $T \approx 6 \times 10^3 \text{ K}$, $B \approx 10^3 \text{ G}$ in corrispondenza delle macchie solari (rappresenta un fluido debolmente ionizzato).
\item \textbf{Nucleo dei reattori Tokamak:} $n_e \approx 10^{14} \text{ cm}^{-3}$, $T \approx 10^8 \text{ K}$, $B \approx 5 \times 10^4 \text{ G}$ (ambiente artificiale creato per lo studio e il mantenimento del confinamento magnetico del plasma).
\item \textbf{Nucleo solare:} $n_e \approx 6 \times 10^{25} \text{ cm}^{-3}$, $T \approx 1.5 \times 10^7 \text{ K}$ con un campo magnetico ignoto o incerto (il plasma interno è totalmente ionizzato e dominato dalla pressione termica radiativa e particellare).
\item \textbf{Interno delle nane bianche:} $n_e \approx 10^{30} \text{ cm}^{-3}$, $T \approx 10^7 \text{ K}$, e campi magnetici capaci di arrivare fino a $10^9 \text{ G}$ (il mezzo si trova in uno stato di materia degenere per via delle altissime densità).
\end{itemize}

Dal momento che un'unica struttura teorica deve avere l'ambizione di poter descrivere tutti questi ambienti simultaneamente, diviene un imperativo metodologico l'utilizzo dei \textbf{parametri adimensionali} per classificare lo stato del mezzo. I parametri di riferimento più importanti per questa disciplina includono il rapporto dimensionale $\lambda_D/L$ (tra lunghezza di Debye e scala del sistema), il parametro di interazione $N_D$, il fattore di accoppiamento termodinamico $\Gamma$, la valutazione del regime di degenerazione quantistica, il tasso collisionale normalizzato $\nu/\omega_p$, e il numero di Reynolds magnetico $R_m$.

\section{Limiti dell'Approccio Deterministico a Singola Particella}

In via di principio teorico, la dinamica di un gas di particelle ionizzate può essere affrontata a livello microscopico tracciando l'evoluzione di ogni singolo elemento. Se si considera un sistema comporto da $N$ cariche, in cui ciascuna è caratterizzata da una propria massa $m_a$ e da una propria carica $q_a$, il moto orbitale deterministico è definito imponendo la seconda legge di Newton in congiunzione con l'espressione formale della forza di Lorentz. In simboli, avremmo un enorme set di equazioni differenziali del secondo ordine:
\begin{equation}
m_a \frac{d^2 \mathbf{r}_a}{dt^2} = q_a \left[ \mathbf{E}(\mathbf{r}_a, t) + \frac{1}{c} \dot{\mathbf{r}}_a \times \mathbf{B}(\mathbf{r}a, t) \right], \quad a = 1, \dots, N
\end{equation}
dove l'indice $a$ scorre su tutte le particelle del sistema. A queste vanno ad accoppiarsi indissolubilmente i campi elettrico $\mathbf{E}$ e magnetico $\mathbf{B}$ governati dalle equazioni di Maxwell. Tali campi, infatti, traggono le loro sorgenti primarie proprio dalle particelle cariche, il cui contributo di densità formale è scrivibile mediante l'introduzione delle distribuzioni impulsive di \textbf{Delta di Dirac} per la densità di carica macroscopica ($\rho_c$) e la densità di corrente ($\mathbf{j}$):
\begin{equation}
\rho_c = \sum{a} q_a \delta(\mathbf{r} - \mathbf{r}*a)
\end{equation}
\begin{equation}
\mathbf{j} = \sum*{a} q_a \dot{\mathbf{r}}_a \delta(\mathbf{r} - \mathbf{r}_a)
\end{equation}

Tuttavia, risolvere analiticamente o numericamente questo sistema per tracciare le particelle risulta fondamentalmente \textbf{impraticabile}. Il numero di cariche contenute perfino in un minuscolo volume di laboratorio è sterminato ($N \sim 10^{23}$), cifra che schizza facilmente fino a ordini di grandezza stellari qualora si studi il Sole ($N \sim 10^{57}$). Questo significherebbe dover gestire simultaneamente $6N$ equazioni differenziali ordinarie (ODE) mutuamente accoppiate e regolate da forze di interazione a lungo raggio qual è la repulsione/attrazione di Coulomb.

Oltre a non essere computazionalmente gestibile, questo formalismo rigoroso non è nemmeno \textbf{desiderabile}. Innanzitutto, non avremo mai la capacità tecnica di conoscere l'enorme spazio delle $6N$ condizioni iniziali necessarie al problema. Inoltre, la formulazione puramente microscopica della legge di Newton non tiene in minima considerazione le interazioni tipiche di un esperimento reale o un ambiente aperto: effetti di confinamento alle pareti, termini sorgente, interazioni dissipative e di scambio quantitativo dovute agli urti (collisioni) con le particelle neutre circostanti, così come lo sviluppo di eventuali reazioni termonucleari, chimiche o di ricombinazione termica risultano del tutto assenti dalle equazioni originali.

Il \textbf{risultato chiave} di questa problematica è l'obbligatoria transizione verso una formulazione profondamente differente, ovvero una descrizione formale di tipo \textbf{statistico}. L'ente matematico che va a sostituire le coordinate spaziotemporali delle singole cariche è la \textbf{funzione di distribuzione cinetica} $f(\mathbf{x}, \mathbf{v}, t)$ (il cui comportamento intrinseco porta all'equazione di Vlasov e costituisce il cuore del capitolo 5 del testo di riferimento). Sarà poi l'integrazione pesata della funzione, attraverso il calcolo matematico dei suoi momenti di vario ordine spaziale, a definire l'avvento formale delle \textbf{variabili fluide} macroscopiche quali la densità numerica $n$, la velocità bulk direzionale $\mathbf{u}$ e la pressione $p$.

\section{Livelli di Descrizione Analitica: La Mappa Concettuale del Corso}

Lo studio teorico e formale della fisica dei plasmi avviene tramite una rigorosa concatenazione di modelli, costruendo una vera e propria mappa concettuale di descrizione su più livelli, in cui ogni "gradino" consiste in una approssimazione specifica matematicamente dominata da un parametro infinitesimo (o molto grande). Questa mappa del corso si divide concettualmente nelle seguenti tappe:

\begin{itemize}
\item \textbf{Descrizione esatta a N-corpi:} La base da cui scaturisce il modello microscopico reale (noto sotto il formalismo rigoroso dell'equazione di Klimontovich).
\item \textbf{Teoria Cinetica:} Facendo tendere il parametro di plasma a valori formidabili (limite in cui $N_D \gg 1$), le perturbazioni granulari del sistema perdono rilevanza e si passa al continuo fluidodinamico, abbracciando la Teoria Cinetica. In quest'area l'evoluzione dipende dallo studio dell'equazione differenziale che descrive la funzione $f(\mathbf{x}, \mathbf{v}, t)$ (argomento dei capitoli dal 5 all'8). Esiste inoltre una ramificazione detta \textbf{teoria delle orbite di singola particella} (sviluppata nei capitoli 3-4), strettamente associata all'approssimazione geometrica del centro di guida ("guiding centre"), la quale connette analiticamente le traiettorie di Lorentz con gli scambi energetici studiati in cinetica.
\item \textbf{Modello a Due Fluidi:} Integrando la Teoria Cinetica sulle distribuzioni e operando un'approssimazione formale nota come chiusura dei momenti (moments closure), il gas ionizzato viene descritto come l'interpenetrazione macroscopica di un fluido elettronico e un fluido ionico (studio presentato nei capitoli 9-10).
\item \textbf{Magnetoidrodinamica (MHD):} Quando si osserva il plasma in regimi in cui le frequenze temporali di sollecitazione sono basse rispetto alla frequenza di ciclotrone dello ione ($\omega \ll \Omega_i$), su dimensioni caratteristiche infinitamente ampie ($L \gg d_i$) tali da poter imporre la validità stretta della quasi-neutralità di carica, il sistema si condensa in un singolo fluido magnetizzato. Le dinamiche derivanti da tali supposizioni fondano la Magnetoidrodinamica (MHD), snocciolata diffusamente nei capitoli 12-15.
\item \textbf{Teoria della Dinamo:} Sotto le condizioni limite in cui gli effetti di convezione magnetica eclissano totalmente i fenomeni diffusivi (limite imposto da un numero di Reynolds magnetico altissimo, $R_m \gg 1$), lo studio magnetoidrodinamico trapassa nella Teoria della Dinamo, fulcro accademico dedicato ai processi di auto-eccitazione ed evoluzione galattico-stellare dei campi (capitoli dal 16 al 21).
\end{itemize}
Come premessa essenziale alla mappa di corso appena esposta, le argomentazioni odierne getteranno luce esclusiva sui parametri propedeutici al primo passo limitante dell'albero formale, definendo i criteri fisici della lunghezza di Debye ($\lambda_D$), le frequenze operative elementari ($\omega_p$), l'impatto topologico espresso dal numero $N_D$, e infine l'incidenza pratica correlata all'intensità o grado di ionizzazione che permea la materia.

\section{Sistema di Unità di Misura: Adozione Formale del Sistema Gaussiano}

Una caratteristica dirimente nell'architettura delle equazioni che formano la Teoria dei Plasmi è la capillare scelta, applicata lungo tutto il corso, di impiegare esclusivamente le convenzioni inerenti al sistema di misura \textbf{Gaussiano (CGS)} al posto del Sistema Internazionale (SI). Scegliendo questo sistema, le equazioni formali base che dettano l'Elettrodinamica vedono affiorare dei prefattori strutturali differenti.

Sotto le convezioni di Gauss, il gruppo delle equazioni di Maxwell assume tale elegante simmetria differenziale:
\begin{equation}
\nabla \cdot \mathbf{E} = 4\pi \rho_c
\end{equation}
\begin{equation}
\nabla \times \mathbf{E} = -\frac{1}{c} \frac{\partial \mathbf{B}}{\partial t}
\end{equation}
\begin{equation}
\nabla \cdot \mathbf{B} = 0
\end{equation}
\begin{equation}
\nabla \times \mathbf{B} = \frac{4\pi}{c} \mathbf{j} + \frac{1}{c} \frac{\partial \mathbf{E}}{\partial t}
\end{equation}
Allo stesso modo, la legge della Forza di Lorentz viene matematicamente emendata ponendo il termine legato alla propagazione della luce $c$ direttamente come coefficiente dimensionale applicato al termine di interazione prodotto vettoriale tra campo cinetico $\mathbf{v}$ e magnetico $\mathbf{B}$:
\begin{equation}
\mathbf{F} = q \left( \mathbf{E} + \frac{\mathbf{v}}{c} \times \mathbf{B} \right)
\end{equation}

Durante l'esposizione della materia risulteranno molto frequenti, nonché essenziali per l'economia strutturale del plasma, alcune composizioni scalari algebriche di tali grandezze vettoriali. Le tre combinazioni più ricorrenti includono:
\begin{itemize}
\item La \textbf{pressione magnetica}, che quantifica la densità di tensione cinetica racchiusa dal campo geomagnetico o stellare ed è computata mediante l'espressione $B^2/8\pi$.
\item La \textbf{velocità di Alfvén}, introdotta come $v_A = B / \sqrt{4\pi \rho}$, grandezza inestimabile per calcolare la propagazione magnetoidrodinamica delle perturbazioni lineari nella densità $\rho$.
\item La \textbf{diffusività magnetica}, misurata dal coefficiente termo-conduttivo $\eta = c^2/4\pi\sigma$, la quale regola lo sfaldamento topologico delle correnti all'interno del continuo conduttore del mezzo.
\end{itemize}

\subsection{Conversione dalle Unità SI ai Parametri Gaussiani e Equivalenze Termodinamiche}
A livello pragmatico è indispensabile familiarizzare con la ricetta per transitare organicamente dalle convenzioni matematiche formulate in Sistema Internazionale alla forma Gaussiana. Più specificamente, l'operazione non si basa su conversioni di magnitudine numerica, ma su \textbf{sostituzioni analitiche degli operatori all'interno dell'algebra delle equazioni stesse}:
\begin{itemize}
\item La permittività dielettrica nel vuoto scala seguendo: $\varepsilon_0 \to 1/4\pi$.
\item La permeabilità magnetica nel vuoto si relaziona mediante la funzione di ritardo relativistico: $\mu_0 \to 4\pi/c^2$.
\item L'induzione del campo magnetico modifica formalmente il proprio accoppiamento dimensionale traslando come: $\mathbf{B} \to \mathbf{B}/c$.
\end{itemize}
Un esempio esplicativo dell'efficacia pragmatica di queste sostituzioni incrociate è dato dal fattore della velocità di Alfvén derivabile prima, in cui l'equazione $B/\sqrt{\mu_0 \rho}$ (presente nello sviluppo fluido sotto la convezione standard SI) diviene meccanicamente, attraverso le equivalenze appena indicate, il termine rigoroso in unità elettrostatiche Gaussiane pari a $B/\sqrt{4\pi \rho}$.

Per una risoluzione esatta dei problemi proposti, è inoltre utile ancorare saldamente il calcolo ad alcuni riferimenti cardinali impiegati nel Sistema:
\begin{itemize}
\item La carica elementare isolata (nel caso dell'elettrone, $e$) equivale al valore quantificato misurabile pari a $4.803 \times 10^{-10} \text{ esu}$.
\item La misura della magnitudine del campo magnetico fa uso del Gauss anziché del Tesla, con l'uguaglianza analitica esatta posta a $1 \text{ G} = 10^{-4} \text{ T}$.
\item L'elettronvolt diviene un viatico fondamentale con la traslazione $1 \text{ eV} = 1.602 \times 10^{-12} \text{ erg}$, parametro quest'ultimo che fissa la corrispondenza termodinamica cinetica al valore asintotico di $11604 \text{ K}$ per ciascun atomo scisso nell'energia di risonanza.
\end{itemize}
Al fine di snellire la formattazione analitica e rimuovere continue digressioni di costanti all'interno dei testi formali, si concorda che le variazioni termiche all'interno delle grandezze dei plasmi vadano espresse \textbf{direttamente in unità di energia}. Questo stratagemma di uniformazione matematica viene portato a compimento annullando artificialmente la costante ponderale standard del modello termodinamico di riferimento, asserendo in modo inequivocabile per tutte le future stesure differenziali l'imposizione costante che $k_B = 1$.

\section{Schermo di Debye e Quasi-neutralità}
Di seguito vengono esaminati i concetti fondamentali relativi al meccanismo dello schermo di Debye e al principio di quasi-neutralità all'interno di un mezzo ionizzato.

\subsection{Una Carica di Prova nel Plasma}
Per comprendere l'effetto di schermo, è utile considerare cosa accade quando si inserisce una carica elettrica in mezzi con proprietà diverse. Se si introduce una carica all'interno di un conduttore perfetto, le cariche libere del materiale si ridistribuiscono istantaneamente fino a quando il campo elettrico interno si annulla del tutto ($E=0$), garantendo così un effetto di schermo perfetto.

In un plasma, tuttavia, la dinamica è differente. Le cariche che compongono il plasma possiedono una notevole energia cinetica ("sono calde"); questo incessante moto di agitazione termica impedisce loro di accumularsi indefinitamente attorno alla carica introdotta. Di conseguenza, nel plasma l'effetto di schermo risulta essere solo \textit{parziale} e si estende su una distanza spaziale finita. Graficamente, se si pone una carica di prova positiva $q > 0$ nell'origine del sistema, si osserva che gli elettroni tendono a raccogliersi attorno ad essa, mentre gli ioni positivi vengono spinti verso l'esterno, confinando l'influenza della carica in una regione di raggio caratteristico $r = \lambda_D$.

Per formalizzare questo fenomeno, consideriamo la carica di prova $q$ fissata nell'origine e un plasma composto da varie specie di particelle, indicate con l'indice $s$. Ciascuna specie è caratterizzata da una carica $q_s$, una densità di particelle imperturbata $n_{s0}$ e una temperatura $T_s$. Assumendo che il plasma di fondo sia complessivamente neutro, la somma delle densità di carica imperturbate deve essere nulla: $\sum_s q_s n_{s0} = 0$. L'equazione di Poisson che governa il potenziale elettrico $\phi$ del sistema deve tenere conto sia della carica puntiforme sorgente sia della distribuzione spaziale delle cariche del plasma:
\begin{equation}
\nabla^2 \phi = -4\pi \left[ q \delta(\mathbf{r}) + \sum_s q_s n_s(\mathbf{r}) \right]
\end{equation}
Per risolvere questa equazione, risulta quindi necessario determinare la funzione di densità spaziale $n_s(\mathbf{r})$ di ciascuna specie in risposta al campo elettrico generato da $q$.

\subsection{La Relazione di Boltzmann dall'Equilibrio delle Forze}
Per ricavare la distribuzione spaziale della densità, partiamo dall'equazione della conservazione della quantità di moto per la specie $s$, trattata secondo il modello a fluido (che verrà approfondito nel capitolo 9):
\begin{equation}
m_s n_s \frac{d\mathbf{u}_s}{dt} = q_s n_s \mathbf{E} - \nabla p_s
\end{equation}
Per semplificare lo studio, imponiamo quattro assunzioni fisiche fondamentali:
\begin{enumerate}
\item Il sistema si trova in un regime \textbf{statico}, il che rende trascurabile il termine di inerzia (il membro di sinistra dell'equazione si annulla).
\item Il campo è puramente \textbf{elettrostatico}, permettendo di esprimere il campo elettrico come il gradiente cambiato di segno del potenziale: $\mathbf{E} = -\nabla\phi$.
\item La temperatura $T_s$ della specie si mantiene \textbf{uniforme} in tutto lo spazio.
\item La specie si comporta come un \textbf{gas ideale}, obbedendo all'equazione di stato $p_s = n_s T_s$.
\end{enumerate}

Sostituendo queste ipotesi nell'equazione di bilancio delle forze, si ottiene:
\begin{equation}
0 = -q_s n_s \nabla\phi - T_s \nabla n_s
\end{equation}
Riorganizzando i termini per separare le variabili, si ricava la forma differenziale:
\begin{equation}
\frac{dn_s}{n_s} = -\frac{q_s}{T_s} d\phi
\end{equation}
Integrando ambo i membri, giungiamo alla relazione:
\begin{equation}
\ln n_s = -\frac{q_s \phi}{T_s} + C
\end{equation}
La costante di integrazione $C$ viene determinata imponendo le condizioni al contorno: a grandissima distanza dalla carica di prova, il potenziale si annulla ($\phi = 0$) e la densità deve tendere al suo valore imperturbato di fondo ($n_s = n_{s0}$). Questa imposizione conduce al risultato chiave, noto come \textbf{Relazione di Boltzmann}:
\begin{equation}
n_s(\mathbf{r}) = n_{s0} \exp\left(-\frac{q_s \phi(\mathbf{r})}{T_s}\right)
\end{equation}

\subsection{Linearizzazione e Equazione di Poisson Schermata}
Sostituendo la densità ottenuta dalla relazione di Boltzmann all'interno dell'equazione di Poisson, costruiamo un'equazione auto-consistente per il potenziale $\phi$:
\begin{equation}
\nabla^2 \phi = -4\pi q \delta(\mathbf{r}) - 4\pi \sum_s q_s n_{s0} e^{-q_s \phi / T_s}
\end{equation}
Questa equazione risulta fortemente non lineare rispetto a $\phi$. Per procedere analiticamente, possiamo applicare una \textbf{linearizzazione}, assumendo di trovarci sufficientemente lontani dalla carica di prova affinché l'energia potenziale elettrostatica sia molto minore dell'energia termica, ovvero $\vert{}q_s \phi\vert{} \ll T_s$. In questo limite, possiamo espandere l'esponenziale arrestandoci al primo ordine ($e^{-y} \simeq 1 - y$). La sommatoria delle densità di carica si approssima quindi a:
\begin{equation}
\sum_s q_s n_s \simeq \sum_s q_s n_{s0} - \phi \sum_s \frac{n_{s0} q_s^2}{T_s}
\end{equation}
Il primo termine a destra dell'uguaglianza ($\sum_s q_s n_{s0}$) rappresenta la carica di fondo del plasma imperturbato ed è identicamente nullo per l'ipotesi di quasi-neutralità. Reinserendo il termine rimanente nell'equazione differenziale originale, si ottiene il risultato fondamentale noto come \textbf{equazione di Poisson schermata}:
\begin{equation}
\nabla^2 \phi - \frac{\phi}{\lambda_D^2} = -4\pi q \delta(\mathbf{r})
\end{equation}
In questa formulazione, il parametro $\lambda_D$ racchiude le proprietà schermanti del mezzo ed è definito mediante il suo inverso al quadrato:
\begin{equation}
\frac{1}{\lambda_D^2} = \sum_s \frac{4\pi n_{s0} q_s^2}{T_s} = \sum_s \frac{1}{\lambda_{Ds}^2}
\end{equation}
Questo dimostra che i contributi di schermo delle singole specie si compongono addizionando i reciproci dei quadrati delle rispettive lunghezze.

\subsection{La Lunghezza di Debye}
Il risultato chiave per quantificare lo schermo è la \textbf{lunghezza di Debye elettronica} ($\lambda_{De}$), la cui espressione in unità Gaussiane è:
\begin{equation}
\lambda_{De} = \left( \frac{T_e}{4\pi n_e e^2} \right)^{1/2} = 743 \left( \frac{T_e / \text{eV}}{n_e / \text{cm}^{-3}} \right)^{1/2} \text{ cm}
\end{equation}
Nel Sistema Internazionale (SI), la medesima grandezza assume la forma $\lambda_{De} = (\varepsilon_0 T_e / n_e e^2)^{1/2}$.

L'analisi di questa formula rivela diverse proprietà fisiche fondamentali legate allo schermo nel plasma:
\begin{itemize}
\item Un plasma più caldo (maggiore $T_e$) produce uno schermo più debole. A temperature elevate, la grande agitazione termica fa sì che le particelle oppongano maggiore resistenza alla compressione necessaria per ammassarsi nella nuvola di schermo.
\item Un plasma più denso (maggiore $n_e$) genera uno schermo più forte, in quanto vi è un numero superiore di cariche disponibili nel volume per neutralizzare la perturbazione.
\item I contributi delle diverse specie si sommano "in parallelo", seguendo la regola $\lambda_D^{-2} = \lambda_{De}^{-2} + \lambda_{Di}^{-2}$. Di conseguenza, la specie caratterizzata dalla lunghezza di Debye minore domina l'effetto di schermo globale.
\item A causa della loro elevata massa, gli ioni sono in grado di schermare unicamente perturbazioni statiche o fenomeni lenti (ovvero dove le velocità in gioco soddisfano $v \ll v_{ti}$). Quando il plasma è interessato da fenomeni molto rapidi, l'inerzia ionica impedisce una ridistribuzione tempestiva, e la lunghezza di Debye efficace del sistema tende a coincidere con la sola lunghezza di Debye elettronica ($\lambda_D \to \lambda_{De}$).
\end{itemize}

\section{Schermo di Debye e Quasi-neutralità}

Nelle sezioni precedenti si è introdotto qualitativamente il concetto di schermo elettrostatico nei plasmi. Ora si procederà con la derivazione analitica rigorosa del potenziale elettrico all'interno del mezzo e con l'analisi delle conseguenze energetiche legate alla formazione della nuvola di schermo attorno alle cariche.

\subsection{Risoluzione dell'Equazione di Poisson Schermata}

L'obiettivo è risolvere l'equazione di Poisson schermata, ricavata linearizzando la relazione di Boltzmann per le densità di carica. Assumendo che il sistema possieda una \textbf{simmetria sferica} attorno alla carica di prova $q$ posta nell'origine, è conveniente cercare una soluzione che si presenti come il prodotto tra il potenziale di Coulomb nudo e una funzione di correzione spaziale $f(r)$. Poniamo dunque l'ansatz:
\begin{equation}
\phi(r) = q \frac{f(r)}{r}
\end{equation}
Per inserire questa espressione nell'equazione di Poisson, dobbiamo calcolarne il laplaciano, $\nabla^2\phi$. Applicando la regola di Leibniz per il laplaciano di un prodotto, $\nabla^2(gh) = g\nabla^2h + 2\nabla g \cdot \nabla h + h\nabla^2g$, e identificando $g = 1/r$ e $h = f$, otteniamo:
\begin{equation}
\nabla^2\left(\frac{f}{r}\right) = \frac{1}{r}\nabla^2f + 2\nabla\left(\frac{1}{r}\right) \cdot \nabla f + f \nabla^2\left(\frac{1}{r}\right)
\end{equation}
Analizziamo questa espressione termine per termine, ricordando che in coordinate sferiche l'operatore laplaciano per funzioni radiali è $\nabla^2 = r^{-2}\partial_r(r^2\partial_r)$ e sfruttando l'identità fondamentale per la singolarità coulombiana $\nabla^2(1/r) = -4\pi\delta(\mathbf{r})$:
\begin{itemize}
\item Il primo termine si espande come: $\frac{1}{r}\nabla^2f = \frac{f''}{r} + \frac{2f'}{r^2}$.
\item Il secondo termine, calcolando i gradienti, diviene: $2\nabla\left(\frac{1}{r}\right) \cdot \nabla f = 2\left(-\frac{\hat{\mathbf{r}}}{r^2}\right) \cdot (f'\hat{\mathbf{r}}) = -\frac{2f'}{r^2}$.
\item Il terzo termine introduce la distribuzione di Dirac: $f \nabla^2\left(\frac{1}{r}\right) = -4\pi f(\mathbf{r}) \delta(\mathbf{r})$.
\end{itemize}
Sommando i contributi, i termini contenenti la derivata prima $f'$ si elidono perfettamente. Inoltre, per le proprietà della funzione delta di Dirac, il prodotto $f(\mathbf{r})\delta(\mathbf{r})$ è non nullo solo nell'origine, potendosi quindi riscrivere come $f(0)\delta(\mathbf{r})$. Il risultato chiave per l'operatore laplaciano applicato alla nostra funzione è dunque:
\begin{equation}
\nabla^2\left(\frac{f}{r}\right) = \frac{f''}{r} - 4\pi f(0) \delta(\mathbf{r})
\end{equation}

Sostituiamo ora questo risultato, unitamente all'ansatz $\phi = qf/r$, all'interno dell'equazione di Poisson schermata completa, $\nabla^2\phi - \phi/\lambda_D^2 = -4\pi q \delta(\mathbf{r})$. Moltiplicando per $q$ si ha:
\begin{equation}
\frac{qf''}{r} - 4\pi q f(0) \delta(\mathbf{r}) - \frac{qf}{r\lambda_D^2} = -4\pi q \delta(\mathbf{r})
\end{equation}
Affinché questa equazione differenziale sia valida in tutto lo spazio, è necessario imporre l'uguaglianza separatamente per le singolarità (i termini con la funzione delta) e per la parte omogenea.
Per quanto riguarda le funzioni delta, i loro coefficienti devono coincidere: $-4\pi q f(0) = -4\pi q$, il che impone la condizione al contorno $f(0) = 1$. Fisicamente, questo assicura che in prossimità della carica di prova ($r \to 0$) il potenziale si comporti esattamente come un potenziale di Coulomb "nudo".

Per distanze $r > 0$, i termini con la delta si annullano, lasciando l'equazione differenziale ordinaria:
\begin{equation}
f'' = \frac{f}{\lambda_D^2}
\end{equation}
La cui soluzione generale è una combinazione lineare di esponenziali:
\begin{equation}
f(r) = A e^{-r/\lambda_D} + B e^{r/\lambda_D}
\end{equation}
Per determinare le costanti, imponiamo le condizioni fisiche al contorno. Il potenziale deve rimanere limitato per $r \to \infty$, il che obbliga a porre $B = 0$ (per sopprimere l'esponenziale divergente). Utilizzando la condizione nell'origine precedentemente trovata, $f(0) = 1$, ricaviamo immediatamente $A = 1$.
Inserendo la funzione $f(r)$ così determinata nel nostro ansatz iniziale, giungiamo al \textbf{Potenziale di Debye-Hückel}:
\begin{equation}
\phi(r) = \frac{q}{r} e^{-r/\lambda_D}
\end{equation}

\subsection{Il Potenziale di Debye e la Funzione Dielettrica del Plasma}

Dal punto di vista spaziale, il potenziale di Debye diverge profondamente dal classico potenziale di Coulomb. Mentre il potenziale coulombiano decade lentamente come $1/r$, il potenziale nel plasma presenta un \textbf{decadimento esponenziale} causato dal fattore $e^{-r/\lambda_D}$. Oltre la distanza caratteristica $r = \lambda_D$, l'influenza della carica crolla drasticamente.

Una prospettiva fisicamente molto istruttiva si ottiene analizzando il problema nello \textbf{spazio di Fourier}. Trasformando l'equazione di Poisson schermata e sostituendo l'operatore laplaciano con $-k^2$, otteniamo:
\begin{equation}
-k^2 \phi_k - \frac{\phi_k}{\lambda_D^2} = -4\pi q
\end{equation}
Risolvendo per la componente di Fourier del potenziale, $\phi_k$, possiamo riscrivere l'espressione mettendo in evidenza $k^2$ a denominatore:
\begin{equation}
\phi_k = \frac{4\pi q}{k^2 + \lambda_D^{-2}} = \frac{4\pi q}{k^2 \epsilon(k)}
\end{equation}
In questa formulazione emerge spontaneamente la \textbf{funzione dielettrica del plasma}, definita come:
\begin{equation}
\epsilon(k) = 1 + \frac{1}{k^2\lambda_D^2}
\end{equation}
Questo risultato dimostra che il plasma si comporta macroscopica come un mezzo dielettrico disperdente. Per scale spaziali ampie, corrispondenti a numeri d'onda piccoli ($k\lambda_D \lesssim 1$), il mezzo interviene schermando pesantemente il potenziale elettrico. Al contrario, esplorando scale spaziali piccolissime, corrispondenti a numeri d'onda elevati ($k\lambda_D \gg 1$), il termine aggiuntivo nella funzione dielettrica diventa trascurabile ($\epsilon(k) \to 1$) e la carica appare "nuda", ovvero percepibile nella sua interezza. Invertendo la trasformata di Fourier di questa espressione si riottene rigorosamente il potenziale $q e^{-r/\lambda_D}/r$ nello spazio reale.

\subsection{La Nuvola di Schermo e la sua Energia}

Per confermare la bontà del modello, è necessario quantificare la carica totale accumulata attorno alla carica di prova. Utilizzando l'approssimazione linearizzata, la \textbf{densità di carica indotta} nel plasma è proporzionale al potenziale stesso:
\begin{equation}
\rho_{\text{ind}} = \sum_s q_s(n_s - n_{s0}) = -\frac{\phi}{4\pi\lambda_D^2}
\end{equation}
Integrando questa densità indotta su tutto il volume tridimensionale, in coordinate sferiche, otteniamo la carica totale della nuvola di schermo:
\begin{equation}
\int \rho_{\text{ind}} d^3r = -\frac{1}{4\pi\lambda_D^2} \int_0^\infty \frac{q}{r} e^{-r/\lambda_D} 4\pi r^2 dr = -\frac{q}{\lambda_D^2} \int_0^\infty r e^{-r/\lambda_D} dr
\end{equation}
L'integrale notevole $\int_0^\infty r e^{-r/\lambda_D} dr$ restituisce esattamente il valore $\lambda_D^2$. Il risultato finale dell'operazione è pertanto \textbf{$-q$}. Ciò dimostra matematicamente che la nuvola di schermo neutralizza in modo esatto e completo la carica di prova introdotta nel mezzo.

A questo punto possiamo determinare l'energia di interazione. Il potenziale totale $\phi$ può essere scomposto nella somma del potenziale di Coulomb nudo e del potenziale generato esclusivamente dalla nuvola di schermo, $\phi_{\text{scr}}$:
\begin{equation}
\phi = \frac{q}{r} + \phi_{\text{scr}}
\end{equation}
Per trovare il potenziale di schermo agente sulla particella stessa, dobbiamo valutare $\phi_{\text{scr}}$ nell'origine ($r \to 0$). Operiamo raccogliendo $q/r$ ed espandendo l'esponenziale in serie di Taylor al primo ordine:
\begin{equation}
\phi_{\text{scr}}(0) = \lim_{r \to 0} \frac{q}{r} \left( e^{-r/\lambda_D} - 1 \right) = \lim_{r \to 0} \frac{q}{r} \left( -\frac{r}{\lambda_D} + \dots \right) = -\frac{q}{\lambda_D}
\end{equation}
Si deduce che ogni particella interagisce con la propria nuvola di schermo acquisendo un'energia potenziale pari a:
\begin{equation}
q\phi_{\text{scr}}(0) = -\frac{q^2}{\lambda_D}
\end{equation}

\subsection{Energia di Correlazione e Parametro di Plasma}

L'analisi svolta per una singola particella di prova si estende naturalmente a tutto il plasma, poiché \textit{ogni} particella del gas ionizzato porta con sé la propria nuvola di schermo. L'\textbf{energia di correlazione} totale per unità di volume, $U_{\text{corr}}/V$, si ottiene sommando l'energia di schermo di tutte le particelle. È tuttavia fondamentale introdurre un fattore $1/2$ per evitare di contare due volte l'interazione tra ogni singola coppia di cariche:
\begin{equation}
\frac{U_{\text{corr}}}{V} = \frac{1}{2} \sum_s n_{s0} q_s \phi_{\text{scr},s}(0) = -\frac{1}{2\lambda_D} \sum_s n_{s0} q_s^2
\end{equation}

Per rendere esplicita l'entità di questa energia, applichiamo la formula al caso di un \textbf{plasma di idrogeno} puro in equilibrio termico, dove le temperature e le densità di elettroni e ioni (protoni) sono uguali: $T_e = T_i = T$ e $n_e = n_i = n$. In queste condizioni, l'inverso del quadrato della lunghezza di Debye complessiva si semplifica in $\lambda_D^{-2} = 8\pi n e^2 / T$.
L'energia di correlazione media per singola particella, indicata con $w$, si calcola dividendo la densità di energia per la densità totale di particelle ($2n$):
\begin{equation}
w = \frac{U_{\text{corr}}}{2nV} = -\frac{e^2}{2\lambda_D}
\end{equation}

Il grado di idealità del plasma è stabilito dal confronto tra il modulo di questa energia di correlazione $\vert{}w\vert{}$ e l'energia cinetica termica caratteristica, $T$. Calcoliamo il loro rapporto, sfruttando l'espressione di $\lambda_D^{-2}$ ricavata per l'idrogeno:
\begin{equation}
\frac{|w|}{T} = \frac{e^2}{2\lambda_D T} = \frac{1}{2\lambda_D} \cdot \left( \frac{1}{8\pi n \lambda_D^2} \right) = \frac{1}{16\pi n \lambda_D^3}
\end{equation}
Riconosciamo nel denominatore il \textbf{parametro di plasma} $N_D$, definito come il numero di particelle contenute in una sfera di Debye: $N_D \equiv \frac{4\pi}{3} n \lambda_D^3$. Sostituendo questa definizione nell'equazione, otteniamo il \textbf{risultato chiave} della termodinamica dei plasmi:
\begin{equation}
\frac{|w|}{T} = \frac{1}{12 N_D}
\end{equation}
Questa relazione sancisce che la condizione affinché il plasma si comporti macroscopicamente come un \textbf{gas ideale} ($\vert{}w\vert{} \ll T$, ovvero l'energia di interazione elettrostatica risulta trascurabile rispetto all'agitazione termica) è matematicamente equivalente a richiedere che vi siano \textbf{molte particelle all'interno di una sfera di Debye} ($N_D \gg 1$).

\section{La Quasi-neutralità nel Plasma}

Per comprendere appieno la natura collettiva del plasma, è fondamentale analizzare il meccanismo e le implicazioni della quasi-neutralità macroscopica. Supponiamo che i moti di agitazione termica all'interno del mezzo inducano una separazione locale delle cariche elettriche, generando uno squilibrio di densità $\delta n = n_i - n_e$ su una data scala spaziale caratteristica $l$. Questa disomogeneità nella distribuzione delle cariche produrrà inevitabilmente una fluttuazione del potenziale elettrico $\delta \phi$, la cui magnitudine può essere stimata come:
\begin{equation}
\delta \phi \sim 4\pi e \delta n l^2
\end{equation}

Questo processo di separazione delle cariche non può proseguire indefinitamente: esso viene arrestato dalle forze di richiamo elettrostatiche non appena l'energia potenziale associata allo squilibrio eguaglia l'energia termica disponibile nel mezzo, ovvero quando si verifica la condizione $e \delta \phi \sim T$. Inserendo questa limitazione nell'espressione precedente, si ottiene un \textbf{risultato chiave} per definire il grado di neutralità del sistema:
\begin{equation}
\frac{\delta n}{n} \sim \frac{T}{4\pi n e^2 l^2} = \frac{\lambda_D^2}{l^2} \ll 1 \quad \text{per} \quad l \gg \lambda_D
\end{equation}

Da questa fondamentale relazione scaturiscono diverse conseguenze fisiche dirette:
\begin{itemize}
\item Su scale spaziali ampiamente superiori alla lunghezza di Debye ($l \gg \lambda_D$), la deviazione relativa della densità di carica ($\delta n / n$) risulta del tutto trascurabile. Pertanto, su scale macroscopiche, si può assumere con un'elevatissima accuratezza che la densità elettronica compensi esattamente la densità di carica ionica totale, $n_e = \sum_i Z_i n_i$.
\item È tuttavia cruciale sottolineare che, sebbene valga l'equivalenza di densità macroscopica, il campo elettrico macroscopico non si annulla rigorosamente ($E \neq 0$). In questo regime, il principio di quasi-neutralità va di fatto a sostituire l'impiego diretto dell'equazione di Poisson. Il campo elettrico $\mathbf{E}$ non viene più determinato calcolando la divergenza $\nabla \cdot \mathbf{E} = 4\pi \rho_c$, bensì viene fissato dall'equazione per il bilancio delle forze sugli elettroni (un concetto che sarà alla base della legge di Ohm generalizzata, esplorata nel capitolo 9).
\item Come corollario logico, un gas ionizzato le cui dimensioni fisiche $L$ risultino minori o confrontabili con lo spessore di schermo termico ($L \lesssim \lambda_D$) si configura come una mera collezione di cariche interagenti a coppie, e non rispetta i criteri per essere definito e trattato come un vero plasma.
\end{itemize}

\section{Limiti della Linearizzazione del Potenziale}

Nella trattazione dell'equazione di Poisson-Boltzmann per il calcolo dello schermo di Debye, è stata impiegata una \textbf{linearizzazione}. Tale semplificazione matematica poggia sull'assunzione che l'energia potenziale elettrostatica sia strettamente inferiore all'energia termica, ossia $e\phi \ll T$. Questo equivale a richiedere che la distanza $r$ dalla carica perturbante sia molto maggiore di un parametro critico $b_0$ ($r \gg b_0$). Questo parametro, noto come \textbf{distanza classica di massimo avvicinamento} (classical distance of closest approach), è definito dalla seguente equazione:
\begin{equation}
b_0 = \frac{e^2}{T}
\end{equation}

Se valutiamo il rapporto tra questa distanza di collisione critica e la lunghezza di Debye, facciamo emergere un diretto collegamento con il parametro di plasma $N_D$:
\begin{equation}
\frac{b_0}{\lambda_D} = \frac{1}{4\pi n \lambda_D^3} = \frac{1}{3 N_D}
\end{equation}

Capire dove e perché la linearizzazione fallisce è cruciale per la solidità della teoria, e ci porta a trarre le seguenti deduzioni:
\begin{itemize}
\item Il "nucleo" non lineare attorno a una singola carica, ovvero la regione ristretta definita da $r \lesssim b_0$, contiene in media un numero di particelle trascurabile, proporzionale al volume $\sim n b_0^3 \sim N_D^{-2} \ll 1$. Dunque, le violazioni della linearità coinvolgono statisticamente meno di una particella, non invalidando la descrizione globale.
\item Trattare la densità di carica attraverso un campo continuo (approccio di \textit{mean field}) necessita della presenza di innumerevoli particelle all'interno della nuvola di Debye, riaffermando come fondamento del plasma la condizione $N_D \gg 1$.
\item Il rapporto adimensionale $\lambda_D/b_0$ riapparirà nella formulazione analitica degli urti col nome di \textbf{Logaritmo di Coulomb} ($\ln \Lambda$). In quel contesto (trattato nel modulo 2a), l'argomento del logaritmo sarà definito come $\Lambda = 12\pi n \lambda_{De}^3$.
\end{itemize}

\subsection{Le Guaine (Sheaths) Non Neutrali}
Un contesto fisico in cui la linearizzazione crolla irrimediabilmente è la regione posta nelle immediate vicinanze di una parete fisica del contenitore o di una sonda materiale intrusiva. In prossimità di questi confini, l'energia potenziale elettrostatica scala ai medesimi ordini di grandezza dell'energia termica ($e\phi \sim T$). Quando ciò accade, l'ipotesi di linearizzazione decade, favorendo la formazione spontanea di una \textbf{guaina} (o \textit{sheath} di Langmuir) puramente non neutrale, il cui spessore ammonta tipicamente ad alcune lunghezze di Debye. Questo è, peraltro, lo stesso esatto meccanismo con cui i grani di polvere solida (\textit{dust grains}) presenti nel mezzo interstellare (ISM) acquisiscono una forte carica elettrica.

\section{Stime della Lunghezza di Debye in Ambienti Astrofisici}

Per concretizzare le differenze di scala dei mezzi fisici, riportiamo le stime caratteristiche calcolate per la lunghezza di Debye ($\lambda_{De}$) e per il parametro di plasma ($N_D$) nei maggiori contesti universali. Tali stime sono ottenute impiegando le espressioni operative in unità Gaussiane:
\begin{equation}
\lambda_{De} = 743 \left(\frac{T_{\text{eV}}}{n}\right)^{1/2} \text{ cm} \quad \text{e} \quad N_D = \frac{4\pi}{3} n \lambda_{De}^3
\end{equation}

Analizzando la tabella dei valori caratteristici per diversi regimi, possiamo tracciare una mappatura quantitativa:
\begin{itemize}
\item \textbf{Corona solare}: Con una densità numerica di elettroni $n_e = 10^9 \text{ cm}^{-3}$ e una temperatura $T = 10^6 \text{ K}$ (corrispondente a $86 \text{ eV}$), il meccanismo di schermo agisce a cortissimo raggio, con $\lambda_{De} = 0.2 \text{ cm}$, includendo comunque all'interno del volume $N_D = 4 \times 10^7$ particelle.
\item \textbf{Mezzo interstellare ionizzato caldo (Warm ionised ISM)}: Estremamente rarefatto ma termicamente vivace ($n_e = 10^{-1} \text{ cm}^{-3}$, $T = 10^4 \text{ K} \sim 0.86 \text{ eV}$), questo gas sfoggia una lunghezza di schermo di ben $20 \text{ m}$, che accoglie $N_D = 4 \times 10^9$ corpuscoli schermanti.
\item \textbf{Gas intra-ammasso (Intracluster gas)}: Date le condizioni estreme di vuoto ($n_e = 10^{-3} \text{ cm}^{-3}$) combinate a energie parossistiche ($T = 10^8 \text{ K} \sim 8.6 \text{ keV}$), l'impronta di schermo raggiunge i $20 \text{ km}$, contenendo un astronomico parametro $N_D = 4 \times 10^{16}$.
\item \textbf{Reattore Tokamak}: Costruito per il confinamento magnetico artificiale, ha plasma nucleare denso e caldissimo ($n_e = 10^{14} \text{ cm}^{-3}$, $T = 10^8 \text{K}$), con un minuscolo raggio d'azione schermante pari a 70 $\mu$ m (dove si contano $N_D = 10^8$ particelle).
\item \textbf{Nucleo solare (Solar core)}: Domina per densità ($n_e = 6 \times 10^{25} \text{ cm}^{-3}$) e calore denso ($T = 1.5 \times 10^7 \text{K}$). Al suo interno, $\lambda_{De}$ collassa a scale atomiche, $3 \times 10^{-9} \text{cm}$, per un bassissimo parametro di plasma approssimato a $N_D \approx 10$.
\end{itemize}

Dalla lettura e comparazione di questi dati emerge che la quasi totalità dei plasmi astrofisici esibisce una natura \textbf{estremamente ideale}, ampiamente sostenuta dall'alta disponibilità di cariche schermanti entro le singole sfere ($N_D \gg 1$).
La peculiare eccezione, di enorme interesse fisico, risiede nel \textbf{nucleo solare}: all'interno del cuore della stella, infatti, la contrazione fa sì che la lunghezza di Debye sia dimensionalmente comparabile alla spaziatura media tra le singole particelle ($n^{-1/3} \approx 2.5 \times 10^{-9} \text{ cm}$). Ne deriva che la correzione teorica di Debye-Hückel applicata alla pressione fluidodinamica porta a deviazioni dell'ordine di $\approx 1\%$. Per quanto contenuta e impercettibile negli altri scenari esposti, questa deviazione pressoria all'interno del centro solare è tangibile e può essere rilevata e tracciata attraverso le sensibili osservazioni proprie dell'eliosismologia.

\section{Oscillazioni di Plasma}

In questa sezione si introduce il concetto di \textbf{oscillazioni di plasma}, analizzando il comportamento collettivo dinamico delle cariche all'interno di un mezzo ionizzato.

\subsection{Il Modello del Foglio Elettronico (Electron Sheet Model)}
Per derivare intuitivamente la frequenza naturale di oscillazione del plasma, è utile ricorrere a un modello unidimensionale semplificato noto come "modello del foglio elettronico". In questo scenario, si ipotizza che gli ioni, essendo molto più massicci, rimangano immobili formando uno sfondo di carica positiva uniformemente distribuita. Supponiamo ora che un intero strato (o "foglio") di elettroni, avente densità imperturbata $n_0$, venga dislocato rigidamente dalla sua posizione di equilibrio di una piccola distanza vettoriale $\xi$.

Questa traslazione macroscopica genera una separazione di carica ai bordi della regione perturbata, venendosi a creare una densità di carica superficiale pari a $\pm e n_0 \xi$. La configurazione spaziale risultante è analoga a quella di un \textbf{condensatore a facce piane e parallele}. Le cariche superficiali inducono un campo elettrico di richiamo uniforme al suo interno, la cui intensità è data da:
\begin{equation}
E = 4\pi e n_0 \xi
\end{equation}

Sotto l'azione di questo campo elettrico auto-generato, ogni singolo elettrone contenuto nel foglio subisce una forza di ripristino proporzionale allo spostamento. Scrivendo l'equazione del moto (seconda legge di Newton) per il singolo elettrone, si ottiene:
\begin{equation}
m_e \ddot{\xi} = -eE = -4\pi n_0 e^2 \xi
\end{equation}
Questa espressione ha la forma canonica dell'equazione di un \textbf{oscillatore armonico semplice}. Riorganizzando i termini, emerge il \textbf{risultato chiave} di questa trattazione, ovvero la definizione della \textbf{frequenza di plasma elettronica} ($\omega_{pe}$):
\begin{equation}
\omega_{pe} = \left( \frac{4\pi n_e e^2}{m_e} \right)^{1/2} = 5.64 \times 10^4 (n_e / \text{cm}^{-3})^{1/2} \text{ rad s}^{-1}
\end{equation}
(Si noti che l'identificativo $n_0$ è stato sostituito con la notazione standard $n_e$ per la densità elettronica).

Questa frequenza intrinseca possiede due caratteristiche fondamentali all'interno di questo modello di \textbf{plasma freddo}: è totalmente \textbf{indipendente} dall'ampiezza dello spostamento iniziale ($\xi$) ed è altrettanto \textbf{indipendente} dalla lunghezza d'onda della perturbazione.

\subsection{Derivazione tramite il Modello a Fluido}

Il risultato ottenuto con il modello geometrico semplificato può essere rigorosamente derivato impiegando la teoria a fluido. Il metodo procedurale qui illustrato costituisce l'impianto logico standard che verrà \textbf{utilizzato ripetutamente per studiare ogni tipo di onda e instabilità} nel corso della disciplina.

\subsubsection{1. Definizione delle Assunzioni e delle Equazioni di Base}
Il modello poggia sulle seguenti assunzioni fisiche:
\begin{itemize}
\item Gli \textbf{ioni sono immobili} e formano uno sfondo neutro di densità uniforme $n_0$. Questa approssimazione è giustificata dal fatto che la massa ionica è immensamente superiore a quella elettronica ($m_i \gg m_e$), per cui gli elettroni reagiscono alle variazioni di campo in tempi ordini di grandezza inferiori rispetto agli ioni.
\item Il fluido elettronico è considerato \textbf{freddo}, imponendo $T_e = 0$. L'annullamento della temperatura implica l'assenza totale di gradienti di pressione termica nel sistema.
\item L'approssimazione è puramente \textbf{elettrostatica}: il campo magnetico è nullo ($\mathbf{B} = 0$) e il campo elettrico deriva esclusivamente da un potenziale scalare ($\mathbf{E} = -\nabla\phi$).
\end{itemize}

Sotto tali ipotesi, la dinamica è descritta da tre equazioni differenziali accoppiate:
\begin{enumerate}
\item \textbf{Equazione di Continuità} per la conservazione della massa elettronica:
\begin{equation}
\frac{\partial n_e}{\partial t} + \nabla \cdot (n_e \mathbf{u}_e) = 0
\end{equation}
\item \textbf{Equazione del Moto (Quantità di moto)} per il fluido elettronico, priva del termine di pressione:
\begin{equation}
m_e n_e \left[ \frac{\partial \mathbf{u}_e}{\partial t} + (\mathbf{u}_e \cdot \nabla)\mathbf{u}_e \right] = e n_e \nabla\phi
\end{equation}
(dove si è tenuto conto del segno della carica dell'elettrone nell'espressione della forza).
\item \textbf{Equazione di Poisson} per il campo auto-generato dallo squilibrio di densità tra lo sfondo ionico fisso ($n_0$) e la densità elettronica fluttuante ($n_e$):
\begin{equation}
\nabla^2\phi = -4\pi e (n_0 - n_e)
\end{equation}
\end{enumerate}

\subsubsection{2. Equilibrio, Perturbazione e Linearizzazione}
Seguendo il metodo standard:
\begin{enumerate}
\item \textbf{Ricerca dell'Equilibrio:} Individuiamo uno stato stazionario (indicato con il pedice 0) che soddisfi banalmente le equazioni. In questo caso, il fluido è fermo ($\mathbf{u}_{e0} = 0$), la densità elettronica è uniforme e bilancia quella ionica ($n_{e0} = n_0$), e il potenziale è nullo ($\phi_0 = 0$).
\item \textbf{Introduzione della Perturbazione:} Introduciamo delle piccole deviazioni spaziotemporali (indicate con il pedice 1) rispetto allo stato di equilibrio:
\begin{equation}
n_e = n_0 + n_1, \quad \mathbf{u}_e = \mathbf{u}_1, \quad \phi = \phi_1
\end{equation}
assumendo che l'ampiezza della fluttuazione di densità sia infinitesima ($n_1 \ll n_0$).
\item \textbf{Linearizzazione:} Sostituiamo le espressioni perturbate nelle tre equazioni di base. Il passaggio cruciale della linearizzazione consiste nel mantenere esclusivamente i termini lineari (del \textbf{primo ordine}) nelle quantità perturbate, eliminando o trascurando tutti i prodotti incrociati di ordine superiore come $n_1\mathbf{u}_1$ o il termine convettivo non lineare $(\mathbf{u}_1 \cdot \nabla)\mathbf{u}_1$.
\end{enumerate}

\subsubsection{3. Sintesi delle Equazioni Linearizzate}
Applicando la linearizzazione, il sistema si riduce alla seguente forma:
\begin{itemize}
\item (1) \textbf{Equazione di Poisson:} Lo sfondo neutro $n_0$ si cancella, lasciando solo la fluttuazione.
\begin{equation}
\nabla^2\phi_1 = 4\pi e n_1
\end{equation}
\item (2) \textbf{Equazione di Continuità:} Scompare il termine non lineare di divergenza.
\begin{equation}
\frac{\partial n_1}{\partial t} = -n_0 \nabla \cdot \mathbf{u}_1
\end{equation}
\item (3) \textbf{Equazione del Moto:} Si elide il termine convettivo.
\begin{equation}
m_e \frac{\partial \mathbf{u}_1}{\partial t} = e \nabla\phi_1
\end{equation}
\end{itemize}

\subsubsection{4. Ricavo dell'Equazione dell'Oscillatore}
Per fondere queste equazioni in un'unica relazione differenziale per la densità $n_1$, procediamo calcolando la \textbf{derivata temporale parziale ($\partial_t$)} dell'equazione di continuità (2) e, parallelamente, prendendo la \textbf{divergenza ($\nabla\cdot$)} dell'equazione del moto (3).
Dal primo passaggio si ottiene:
\begin{equation}
\frac{\partial^2 n_1}{\partial t^2} = -n_0 \nabla \cdot \frac{\partial \mathbf{u}_1}{\partial t}
\end{equation}
Sostituendo il termine di divergenza della derivata della velocità ricavato dall'equazione (3):
\begin{equation}
-n_0 \nabla \cdot \frac{\partial \mathbf{u}_1}{\partial t} = -\frac{n_0 e}{m_e} \nabla^2\phi_1
\end{equation}
A questo punto, sfruttiamo l'equazione di Poisson (1) per rimpiazzare il laplaciano del potenziale $\nabla^2\phi_1$ con $4\pi e n_1$. Sostituendo tale espressione, l'equazione collassa nella forma finale:
\begin{equation}
\frac{\partial^2 n_1}{\partial t^2} = -\frac{4\pi n_0 e^2}{m_e} n_1
\end{equation}

Questo risultato analitico coincide \textbf{esattamente} con quello dedotto in modo euristico tramite il modello del foglio elettronico. Otteniamo infatti l'equazione di un oscillatore armonico:
\begin{equation}
\frac{\partial^2 n_1}{\partial t^2} + \omega_{pe}^2 n_1 = 0 \quad \text{con} \quad \omega_{pe}^2 = \frac{4\pi n_0 e^2}{m_e}
\end{equation}
La cui soluzione generale è una fluttuazione armonica nel tempo:
\begin{equation}
n_1(\mathbf{r}, t) = A(\mathbf{r}) \cos[\omega_{pe} t + \alpha(\mathbf{r})]
\end{equation}

\subsection{Analisi della Soluzione: Un'Oscillazione, Non Un'Onda Propagantesi}
L'esame della soluzione $n_1(\mathbf{r}, t)$ porta a conclusioni fisiche di primaria importanza sulla natura della frequenza di plasma.

\begin{itemize}
\item La frequenza $\omega_{pe}$ è formalmente identica a quella trovata col modello del foglio, ma in questo contesto a fluido essa è valida per un \textbf{profilo spaziale iniziale della perturbazione} $A(\mathbf{r})$ del tutto \textbf{arbitrario}.
\item Nell'equazione differenziale finale \textbf{non sopravvivono derivate spaziali} che accoppino punti vicini. Ciò significa che ogni singolo punto del mezzo "oscilla per conto suo", salendo e scendendo localmente con un periodo rigorosamente fissato a $2\pi/\omega_{pe}$.
\item Se si tentasse di imporre una soluzione sotto forma di onda piana propagantesi del tipo $e^{i(kx - \omega t)}$, si scoprirebbe che la relazione di dispersione impone $\omega = \omega_{pe}$ per \textit{qualsiasi} valore del numero d'onda $k$. Poiché la frequenza è costante e indipendente da $k$, la \textbf{velocità di gruppo} $d\omega/dk$ (che rappresenta la velocità di propagazione dell'informazione e dell'energia) risulta essere rigorosamente \textbf{zero}. Pertanto, questa oscillazione stazionaria \textbf{non trasporta energia} attraverso lo spazio.
\item Qual è la causa fisica di questa mancanza di propagazione? L'assunzione cruciale è stata imporre un modello di \textbf{plasma freddo} ($T_e = 0$). L'assenza di agitazione termica implica la mancanza di gradienti di pressione; senza pressione, non esiste alcuna forza meccanica o termodinamica in grado di \textbf{accoppiare elementi di fluido adiacenti} per trasferire il moto da una cella spaziale a un'altra. (Come si vedrà, l'introduzione della pressione termica cambierà radicalmente questo comportamento, permettendo all'oscillazione di propagarsi come un'onda di Langmuir).
\item Infine, il medesimo impianto logico e calcolistico può essere applicato agli \textbf{ioni} in un ipotetico regime in cui si suppongano mobili. Ponendo $m_i$ al posto di $m_e$, si definisce una \textbf{frequenza di plasma ionica} $\omega_{pi}$, correlata a quella elettronica dal rapporto delle masse: $\omega_{pi} = \omega_{pe}(Z^2 m_e / m_i)^{1/2}$. Per i protoni (plasma di idrogeno), tale frequenza risulta essere circa un 43esimo di quella elettronica ($\omega_{pi} \approx \omega_{pe}/43$). In un sistema misto in cui entrambe le specie siano considerate mobili in un regime freddo, la frequenza naturale totale combinata risulterebbe $\omega^2 = \omega_{pe}^2 + \omega_{pi}^2$.
\end{itemize}

\section{Correzione Termica: La Relazione di Bohm-Gross}

Nella trattazione precedente del modello a fluido, si era assunto il limite ideale di un plasma freddo. Ora è fondamentale introdurre la correzione termica ripristinando il termine di pressione all'interno dell'equazione del moto per il fluido elettronico. L'equazione linearizzata assume pertanto la forma:
\begin{equation}
m_e n_0 \frac{\partial \mathbf{u}_1}{\partial t} = e n_0 \nabla\phi_1 - \nabla p_1
\end{equation}
. Assumendo che le trasformazioni avvengano in regime adiabatico, la pressione obbedisce alla legge di proporzionalità $p \propto n^\gamma$. Di conseguenza, la perturbazione di pressione linearizzata si esprime come $p_1 = \gamma T_e n_1$. Poiché stiamo modellando una compressione unidimensionale rapida (fast 1D compression) in cui non vi è il tempo per instaurare un flusso di calore sulla scala della lunghezza d'onda, il sistema possiede un unico grado di libertà termodinamico ($f = 1$). Da ciò deriva che il coefficiente adiabatico risulta essere $\gamma = (f+2)/f = 3$.

Ripetendo i passaggi metodologici già applicati nel caso freddo — ovvero calcolando la divergenza dell'equazione della quantità di moto e la derivata temporale $\partial_t$ dell'equazione di continuità — si perviene a un'equazione differenziale aggiornata per la perturbazione di densità $n_1$:
\begin{equation}
\frac{\partial^2 n_1}{\partial t^2} = -\omega_{pe}^2 n_1 + \frac{3 T_e}{m_e} \nabla^2 n_1
\end{equation}
. Ipotizzando una soluzione sotto forma di onda piana proporzionale a $e^{i(kx - \omega t)}$, possiamo operare le consuete sostituzioni nel dominio di Fourier per gli operatori differenziali: $\partial_t \to -i\omega$ e $\nabla^2 \to -k^2$. Questo procedimento algebrico conduce al \textbf{risultato chiave} della trattazione, noto come \textbf{relazione di Bohm-Gross}:
\begin{equation}
\omega^2 = \omega_{pe}^2 + 3 k^2 v_{te}^2
\end{equation}
dove $v_{te} = (T_e / m_e)^{1/2}$ rappresenta la velocità termica degli elettroni.

L'introduzione della pressione termica ha conseguenze fisiche drastiche: la pressione funge da meccanismo di accoppiamento meccanico tra punti spaziali adiacenti. A differenza dell'oscillazione stazionaria tipica del plasma freddo, la perturbazione ora \textbf{si propaga} nello spazio, acquisendo una velocità di gruppo finita definita da $v_g = 3 k v_{te}^2 / \omega$. Un quadro ancora più completo si ottiene applicando la teoria cinetica, la quale aggiungerà a questa propagazione gli effetti dello smorzamento di Landau (\textit{Landau damping}), come verrà esplorato nel capitolo 6.

\section{Scale Temporali e Spaziali della Risposta del Plasma}

L'interazione tra le dinamiche oscillatorie e il meccanismo di schermo elettrostatico porta alla definizione delle scale caratteristiche spaziotemporali del plasma. Un \textbf{risultato chiave} mette in correlazione le tre grandezze fondamentali studiate finora (frequenza di plasma elettronica $\omega_{pe}$, lunghezza di Debye $\lambda_{De}$ e velocità termica degli elettroni $v_{te}$) tramite l'identità:
\begin{equation}
\omega_{pe} \lambda_{De} = v_{te}
\end{equation}
. Sotto il profilo fisico, ciò implica che la quasi-neutralità di carica venga ripristinata in un intervallo di tempo dell'ordine di $\omega_{pe}^{-1}$. Questo lasso di tempo equivale esattamente al tempo necessario affinché un elettrone termico attraversi una distanza pari alla lunghezza di Debye $\lambda_D$.

Tale profonda connessione struttura i diversi regimi di risposta del plasma:
\begin{itemize}
\item Nei fenomeni caratterizzati da frequenze modeste ($\omega \ll \omega_{pe}$) e scale spaziali estese ($L \gg \lambda_D$), gli elettroni riescono a muoversi in modo abbastanza tempestivo da preservare la neutralità macroscopica del plasma. Questo comportamento collettivo e globalmente neutro fa da fondamento al regime descritto dalla Magnetoidrodinamica (MHD).
\item Affinché l'oscillazione collettiva possa manifestarsi e persistere, le frequenti collisioni con le particelle neutre non devono interromperne il ciclo. Pertanto, il prodotto tra la frequenza di plasma e il tempo caratteristico di collisione con i neutri deve soddisfare la disuguaglianza $\omega_{pe} \tau_n \gg 1$.
\item La frequenza di plasma costituisce inoltre una barriera invalicabile per la propagazione radiativa. Le onde elettromagnetiche (EM) che possiedono una frequenza inferiore a quella locale del mezzo ($\omega < \omega_{pe}$) non sono in grado di propagarsi e vengono totalmente riflesse, un fenomeno le cui implicazioni fisiche saranno trattate nel capitolo 11.
\end{itemize}

\section{La Frequenza di Plasma in Natura}

Lungi dall'essere un puro formalismo teorico, la frequenza di plasma costituisce uno strumento fondamentale per sondare fenomeni osservabili nel mondo naturale e decodificare le proprietà degli ambienti in cui si verificano.

\subsection{La Diagnostica della Ionosfera}
Un primo importante ambito applicativo è l'esplorazione remota della \textbf{ionosfera} terrestre. Sfruttando la riflessione delle onde radio (come documentato storicamente dagli ionogrammi ottenuti da impianti quali HAARP), è possibile ricavare preziosi profili di densità. Trasmettendo un segnale verso l'alto, si ottengono echi di ritorno generati proprio all'altitudine in cui la frequenza dell'onda eguaglia la frequenza di plasma locale, soddisfacendo la condizione $f = f_{pe}(h)$. Rilevando la massima frequenza riflessa (associata allo strato F2 e nota come $f_o F_2$), che si aggira intorno a $4.8 \text{ MHz}$, si deduce il picco di densità elettronica dello strato. Invertendo la formula, $n_e = (f / 8.98 \text{ kHz})^2$, si ottiene un picco di densità ionosferica pari a circa $3 \times 10^5 \text{ cm}^{-3}$.

\subsection{Le Raffiche Solari (Type III Bursts)}
Nel contesto astrofisico stellare, il concetto di frequenza di plasma risulta chiave per l'interpretazione delle \textbf{raffiche solari di Tipo III (Type III bursts)}. Tali emissioni si originano quando intensi fasci di elettroni, eiettati a causa dei brillamenti solari (\textit{flares}) a frazioni consistenti della velocità della luce (ad esempio $c/3$), viaggiano verso l'esterno attraverso l'atmosfera coronale. Il loro passaggio turbolento induce il plasma circostante a irradiare energia elettromagnetica alla frequenza di plasma fondamentale $f_{pe}$ (definita F) e alla sua prima armonica $2f_{pe}$ (definita H). Poiché i fasci attraversano rapidamente regioni in espansione caratterizzate da una densità di materia via via minore, la frequenza radiativa registrata mostra una marcata "deriva" (drift), spostandosi verso il basso nello spettro radio nel giro di pochissimi secondi.

\section{Stime delle Frequenze di Plasma nei Diversi Ambienti}

In molti contesti sperimentali e applicativi, è preferibile indicare la frequenza di plasma nella sua forma lineare operativa ($f_{pe}$) anziché nella sua forma angolare ($\omega_{pe}$). La relazione pratica di conversione è data da:
\begin{equation}
f_{pe} = \frac{\omega_{pe}}{2\pi} = 8.98 \text{ kHz} \left( \frac{n_e}{\text{cm}^{-3}} \right)^{1/2}
\end{equation}
. Attraverso questa equazione diretta, è possibile calcolare agilmente gli ordini di grandezza attesi in regimi di plasma di varia natura:
\begin{itemize}
\item Nel \textbf{gas intra-ammasso (Intracluster gas)}, l'estrema rarefazione ($n_e = 10^{-3} \text{ cm}^{-3}$) restituisce oscillazioni bassissime, nell'ordine dei $280 \text{ Hz}$.
\item Il \textbf{vento solare a 1 Unità Astronomica (1 AU)}, con densità media $n_e = 5 \text{ cm}^{-3}$, manifesta una frequenza di plasma vicina ai $20 \text{ kHz}$.
\item Nella \textbf{ionosfera} ($n_e = 10^6 \text{ cm}^{-3}$), la frequenza sale al regime radio di $9 \text{ MHz}$.
\item La \textbf{corona solare} esibisce plasmi più densi ($n_e = 10^9 \text{ cm}^{-3}$), determinando oscillazioni naturali sui $280 \text{ MHz}$.
\item Nei dispositivi per il confinamento magnetico terrestre, come i \textbf{reattori Tokamak} ($n_e = 10^{14} \text{ cm}^{-3}$), l'innalzamento della densità porta la frequenza nello spettro delle microonde a $90 \text{ GHz}$.
\end{itemize}

Questi parametri numerici permettono inferenze scientifiche immediate, come nel caso di stima di un \textbf{Type III burst}. Qualora un radiotelescopio osservi una di queste raffiche solari con un picco a $100 \text{ MHz}$, ci si può legittimamente chiedere quale sia la densità locale del plasma emissivo. Utilizzando la formula inversa, $n_e = (f / 8.98 \text{ kHz})^2$, si perviene a un risultato di $n_e \approx 1.2 \times 10^8 \text{ cm}^{-3}$. Confrontando tale dato con i profili atmosferici solari standard, possiamo identificare come luogo sorgente dell'emissione le regioni di plasma collocate nella \textbf{bassa corona}.

\section{Riepilogo del Modulo 1a e Prospettive Future}

Concludendo questa prima estesa panoramica sulle proprietà fondamentali, è opportuno sistematizzare le scoperte ricavate sulla definizione dello stato di plasma:
\begin{itemize}
\item A livello microscopico e macroscopico, le singole cariche libere risultano schermate all'interno del gas, con un'influenza elettrostatica che svanisce tipicamente oltre una sfera delimitata dalla \textbf{lunghezza di Debye}, $\lambda_D = (T/4\pi n e^2)^{1/2}$.
\item Gli elettroni vincolati dalle forze di richiamo di queste interazioni Coulombiane sono in grado di instaurare oscillazioni collettive la cui pulsazione propria è determinata dalla \textbf{frequenza di plasma}, $\omega_{pe} = (4\pi n e^2/m_e)^{1/2}$. La congiunzione di questi due concetti teorici si raccorda tramite l'evidenza termodinamica $\omega_{pe} \lambda_D = v_{te}$.
\item Perché l'entità complessa che stiamo analizzando possieda la prerogativa fisica di "plasma", si rende necessaria la rigorosa compresenza di tre disuguaglianze strutturali: il dominio spaziale totale deve surclassare ampiamente le nuvole di schermo interne ($L \gg \lambda_D$), le sfere di Debye devono palesare un alto popolamento statistico per originare il comportamento coeso ($N_D \gg 1$), e la reattività temporale collettiva deve sovrastare senza dubbio gli effetti dispersivi delle collisioni neutre ($\omega_{pe}\tau_n \gg 1$).
\end{itemize}

\subsection{Contenuti del Prossimo Modulo (1b)}
L'itinerario logico delle prossime lezioni teoriche approfondirà le condizioni che regolano il mantenimento dell'equilibrio e l'effettivo instaurarsi del gas ionizzato. Specificamente, il Modulo 1b indagherà le forze di \textbf{accoppiamento} molecolare, la meccanica quantistica della \textbf{degenerazione} e la classificazione per mappe termodinamiche tramite il \textbf{diagramma $T-n$ (Temperatura - Densità)}. In parallelo, verrà derivato formalmente il bilancio radiativo e collisionale dei processi di \textbf{ionizzazione}, esplorando i dettagli applicativi della nota \textbf{equazione di Saha}.